\documentclass[12pt, a4paper]{article}

% Paquetes Esenciales
\usepackage[utf8]{inputenc}
\usepackage[T1]{fontenc}
\usepackage[spanish, es-noshorthands]{babel}
\usepackage{amsmath}
\usepackage{geometry}
\usepackage[american]{circuitikz}
\usepackage{pgfplots}
\pgfplotsset{compat=1.18}

% Configuración de la página
\geometry{a4paper, margin=2.5cm}

\title{\textbf{Análisis AC: Ventilador de 3 Velocidades}}
\author{}
\date{}

\begin{document}

\maketitle

\section*{1. Velocidad Alta (Nodos 'a' y 'd' cortocircuitados)}
Al cortocircuitar los nodos $a$ y $d$ del circuito selector de velocidades, estamos conectando directamente la fuente de tensión alterna al nodo $a$. En esta configuración, los terminales $b$ y $c$ quedan en circuito abierto (al aire), por lo que por las bobinas $L_3$ y $L_4$ no circulará corriente.

El circuito equivalente se reduce a una fuente de tensión en paralelo con dos ramas principales que van hacia el neutro (tierra):
\begin{itemize}
    \item \textbf{Rama 1:} Compuesta por el capacitor $C$ en serie con el inductor $L_1$.
    \item \textbf{Rama 2:} Compuesta por el inductor $L_2$.
\end{itemize}

\begin{center}
    \begin{circuitikz}[scale=1.2, transform shape, thick]
        % Raíl inferior (Neutro)
        \draw (-1,0) -- (3,0);
        \node[ground] at (1,0) {};
        
        % Fuente de AC
        \draw (-1,0) to[sV, l=$V_o$, a={60 Hz}] (-1,4.5) -- (1.5,4.5) node[circ, label=above:$a\,(d)$] (a) {};
        \draw (1.5,4.5) -- (3,4.5);
        
        % Rama L2
        \draw (1.5,0) node[circ]{} to[L, l=$L_2$, i<=$I_2$] (1.5,4.5);
        
        % Rama C - L1
        \draw (3,0) to[C, l=$C$] (3,2.25) to[L, l=$L_1$, i<=$I_1$] (3,4.5) node[circ]{};
    \end{circuitikz}
\end{center}

\section*{Análisis de Impedancias}
En el dominio fasorial, considerando la frecuencia angular $\omega = 2\pi f$ (con $f = 60\text{ Hz}$), las impedancias de cada rama son:

\subsection*{Impedancia de la Rama 1 ($Z_1$)}
La rama 1 tiene el capacitor y el inductor $L_1$ en serie:
\begin{equation}
    Z_1 = Z_C + Z_{L1} = \frac{1}{j\omega C} + j\omega L_1 = j\left(\omega L_1 - \frac{1}{\omega C}\right)
\end{equation}

\subsection*{Impedancia de la Rama 2 ($Z_2$)}
La rama 2 consta únicamente del inductor $L_2$:
\begin{equation}
    Z_2 = Z_{L2} = j\omega L_2
\end{equation}

\section*{Cálculo de Corrientes}
Tomando la tensión de la fuente como referencia fasorial $\tilde{V} = V_o \angle 0^\circ$, determinamos las corrientes en cada rama.

La corriente por la rama 1 ($\tilde{I}_1$) y la corriente por la rama 2 ($\tilde{I}_2$) son:
\begin{equation}
    \tilde{I}_1 = \frac{\tilde{V}}{Z_1} = -j \frac{V_o}{\omega L_1 - \frac{1}{\omega C}}, \quad \quad \tilde{I}_2 = \frac{\tilde{V}}{Z_2} = -j \frac{V_o}{\omega L_2}
\end{equation}

\subsection*{Corriente Total ($\tilde{I}_{total}$)}
Por la Ley de Corrientes de Kirchhoff en el nodo superior, la corriente total de entrada es la suma fasorial:
\begin{equation}
    \tilde{I}_{total} = \tilde{I}_1 + \tilde{I}_2 = -j V_o \left[ \frac{1}{\omega L_1 - \frac{1}{\omega C}} + \frac{1}{\omega L_2} \right]
\end{equation}

\section*{Expresiones en el Dominio del Tiempo y Gráfica}
Para obtener las expresiones temporales, aplicamos la transformada fasorial inversa tomando como referencia la tensión $v(t) = V_o \cos(\omega t)$ (es decir, $\tilde{V} = V_o \angle 0^\circ$).

Para la \textbf{Rama 2} (puramente inductiva), la corriente está retrasada $90^\circ$ respecto a la tensión:
\begin{equation}
    i_2(t) = \frac{V_o}{\omega L_2} \cos(\omega t - 90^\circ) = \frac{V_o}{\omega L_2} \sin(\omega t)
\end{equation}

Para la \textbf{Rama 1} (rama de arranque), en un motor monofásico suele diseñarse para que el efecto capacitivo domine sobre el inductivo ($\frac{1}{\omega C} > \omega L_1$). Esto hace que la reactancia equivalente sea negativa y, por tanto, la corriente adelante $90^\circ$ a la tensión:
\begin{equation}
    \tilde{I}_1 = j \frac{V_o}{\frac{1}{\omega C} - \omega L_1} = \frac{V_o}{\frac{1}{\omega C} - \omega L_1} \angle 90^\circ
\end{equation}
Lo cual, en el dominio del tiempo, se convierte en:
\begin{equation}
    i_1(t) = \frac{V_o}{\frac{1}{\omega C} - \omega L_1} \cos(\omega t + 90^\circ) = - \frac{V_o}{\frac{1}{\omega C} - \omega L_1} \sin(\omega t)
\end{equation}

A continuación se muestra la gráfica normalizada de las formas de onda en este modelo ideal (puramente reactivo), ilustrando cómo $i_1(t)$ está en adelanto y $i_2(t)$ en retraso. Nótese que en este modelo sin resistencias, se observa un desfase de exactamente $180^\circ$ entre ambas corrientes de las ramas, mientras que en un motor real con pérdidas resistivas, el desfase se diseña para estar más cercano a los $90^\circ$.

\begin{center}
    \begin{tikzpicture}
        \begin{axis}[
            width=14cm, height=7cm,
            domain=0:2*pi,
            samples=100,
            axis lines=center,
            xlabel={$\omega t$ (rad)},
            ylabel={Amplitud normalizada},
            xtick={0, 1.5708, 3.14159, 4.71238, 6.28318},
            xticklabels={$0$, $\frac{\pi}{2}$, $\pi$, $\frac{3\pi}{2}$, $2\pi$},
            ytick=\empty,
            legend style={at={(0.5,-0.15)}, anchor=north, legend columns=-1},
            enlargelimits=true,
            clip=false
        ]
        \addplot[thick, black] {cos(deg(x))}; \addlegendentry{$v(t)$}
        \addplot[thick, blue] {0.8*sin(deg(x))}; \addlegendentry{$i_2(t)$ (retraso)}
        \addplot[thick, red] {-0.6*sin(deg(x))}; \addlegendentry{$i_1(t)$ (adelanto)}
        \end{axis}
    \end{tikzpicture}
\end{center}

\section*{El Circuito de Arranque ($C$-$L_1$)}
El arreglo formado por el capacitor $C$ y la bobina $L_1$ se denomina ``circuito de arranque'' (o rama auxiliar) debido a un principio físico fundamental en el funcionamiento de los motores de inducción monofásicos. A continuación se detalla su necesidad y funcionamiento:

\subsection*{1. El problema de los motores monofásicos}
Un motor de inducción necesita que en su interior se genere un \textbf{campo magnético rotatorio} para ``arrastrar'' al rotor y hacerlo girar. Sin embargo, si conectáramos la corriente alterna únicamente a la bobina principal de trabajo ($L_2$), el campo magnético generado solo sería \textit{pulsante} (aumentaría y disminuiría en el mismo eje, pero no giraría). Si un ventilador no tuviera el circuito $C$-$L_1$, al encenderlo el motor simplemente vibraría o zumbaría sin girar, ya que su \textbf{torque de arranque es igual a cero}.

\subsection*{¿Qué es exactamente un campo magnético rotatorio? (Una analogía sencilla)}
Para entender esto, imagina una brújula (que representaría el rotor metálico del motor) y un imán en tu mano.
\begin{itemize}
    \item Si acercas el imán a la brújula y lo dejas quieto (o lo acercas y lo alejas en línea recta), la aguja apuntará hacia él y se quedará atrapada temblando. Esto es similar a un \textbf{campo magnético pulsante}: atrae al rotor, pero no lo hace girar continuamente.
    \item Si comienzas a mover tu mano haciendo círculos con el imán alrededor de la brújula, la aguja empezará a dar vueltas tratando de seguir al imán. ¡Eso es un \textbf{campo magnético rotatorio}!
\end{itemize}
En el motor del ventilador \textbf{no hay imanes moviéndose físicamente}. En su lugar, hay bobinas de cobre fijas alrededor del rotor. Al enviarle dos corrientes que alcanzan su máxima fuerza en momentos distintos (gracias al ``retraso/adelanto'' que explicamos con el capacitor), estas bobinas fijas se vuelven electroimanes que se energizan en una secuencia temporal. Funciona igual que las \textbf{luces de una marquesina} que se encienden una tras otra dando una ilusión perfecta de movimiento; aquí, la fuerza magnética se desplaza en círculo y el rotor gira para perseguirla.

\subsection*{2. La solución: crear una segunda fase artificial}
Para lograr un campo magnético que rote, necesitamos aplicar al menos dos corrientes que estén desfasadas en el tiempo, y que pasen por bobinas ubicadas en posiciones físicas diferentes dentro del motor. El arreglo $C$-$L_1$ sirve exactamente para crear esta segunda fase.

\subsection*{3. El papel clave del capacitor ($C$)}
Como se observó en el análisis fasorial y temporal:
\begin{itemize}
    \item En la rama principal ($L_2$), por ser pura inductancia, la corriente \textbf{se retrasa} respecto al voltaje.
    \item En la rama de arranque ($C$ en serie con $L_1$), el capacitor se elige para que su reactancia capacitiva cancele y supere a la reactancia inductiva de $L_1$. Esto provoca que la corriente en esta rama \textbf{se adelante} respecto al voltaje.
\end{itemize}

\subsection*{4. El empuje inicial (Torque de arranque)}
Gracias a este arreglo, tenemos dos corrientes cruzando por el motor al mismo tiempo: una atrasada y otra adelantada. Esta diferencia temporal de fases (idealmente cercana a $90^\circ$), sumada a la disposición física de las bobinas $L_1$ y $L_2$ en el estator, genera el deseado campo magnético rotatorio. Este campo es el que da el ``empujón'' inicial para que las aspas comiencen a girar.

\subsection*{5. La secuencia temporal y espacial (``Completando la vuelta'')}
La intuición de que una corriente ``completa la vuelta que no pudo terminar la otra'' es conceptualmente excelente, pero con un matiz importante: no es que la corriente viaje completando un círculo de manera eléctrica, sino que son sus \textbf{efectos magnéticos} los que completan el círculo en el espacio físico del motor.

Imaginemos cómo interactúan el tiempo y el espacio:
\begin{itemize}
    \item \textbf{En el espacio:} Las bobinas de la rama de arranque ($L_1$) y de la rama principal ($L_2$) están enrolladas en posiciones físicas contiguas pero desplazadas alrededor del estator.
    \item \textbf{En el tiempo:} Como $i_1(t)$ está adelantada temporalmente, alcanza su pico de fuerza magnética \textit{antes} que $i_2(t)$.
    \item \textbf{El giro:} Primero ocurre un fuerte tirón magnético en la posición de la bobina $L_1$. Cuando este empieza a decaer, la corriente $i_2(t)$ alcanza su pico, provocando un nuevo tirón magnético en la posición contigua correspondiente a la bobina $L_2$. Para el rotor en el centro, esto se siente como una fuerza que se desplaza de una bobina a la siguiente; en efecto, el campo de la segunda bobina toma el relevo y continúa el trabajo que empezó la primera, logrando la rotación ininterrumpida.
\end{itemize}

\vspace{0.5cm}
\noindent\textbf{Nota:} En la mayoría de los ventiladores de techo y de pedestal, se usa un motor de \textit{capacitor de división permanente} (Permanent-Split Capacitor). En este diseño, la rama de arranque $C$-$L_1$ nunca se desconecta mediante un interruptor centrífugo (como sí ocurre en bombas o lavadoras), ya que además de proporcionar el torque de arranque, ayuda a mantener un giro más suave y silencioso durante la operación normal del ventilador.

\newpage
\section*{2. Efecto de la Bobina $L_3$: Velocidad Media (Nodos 'b' y 'd' cortocircuitados)}

Para analizar la siguiente configuración del ventilador, cortocircuitamos los nodos $b$ y $d$. En este escenario, la fuente de tensión alterna se conecta al nodo $b$. El terminal $c$ queda en circuito abierto, por lo que la bobina $L_4$ no tiene participación. Sin embargo, la corriente ahora debe atravesar la bobina $L_3$ para poder llegar al nodo $a$ y alimentar el resto del motor (las ramas de $L_1$ y $L_2$).

\begin{center}
    \begin{circuitikz}[scale=1.2, transform shape, thick]
        % Raíl inferior (Neutro)
        \draw (-1,0) -- (4,0);
        \node[ground] at (1,0) {};
        
        % Fuente de AC conectada a b
        \draw (-1,0) to[sV, l=$V_o$, a={60 Hz}] (-1,4.5) -- (-0.2,4.5) node[circ, label=above:$b\,(d)$] (b) {};
        
        % Inductor L3 actuando de impedancia serie
        \draw (-0.2,4.5) to[L, l=$L_3$, i={\hspace{16pt}$I_{total}$}] (1.8,4.5) -- (2.5,4.5) node[circ, label=above:$a$] (a) {};
        
        % Rama L2
        \draw (2.5,0) node[circ]{} to[L, l=$L_2$, i<=$I_2$] (2.5,4.5);
        
        % Rama C - L1
        \draw (4,0) to[C, l=$C$] (4,2.25) to[L, l=$L_1$, i<=$I_1$] (4,4.5) -- (a);
        \draw (2.5,0) -- (4,0);
    \end{circuitikz}
\end{center}

\subsection*{El Efecto de $L_3$ (Divisor de Tensión e Impedancia Serie)}
Al introducir $L_3$ en el trayecto principal, la bobina actúa como una \textbf{impedancia en serie} con el arreglo en paralelo del motor (compuesto por la rama principal y la rama de arranque). 

Si definimos la impedancia equivalente del motor en el nodo $a$ como $Z_a = Z_1 \parallel Z_2$, la nueva impedancia total vista desde la fuente será:
\begin{equation}
    Z_{total\_media} = Z_{L3} + Z_a = j\omega L_3 + \left( \frac{Z_1 Z_2}{Z_1 + Z_2} \right)
\end{equation}

Este aumento en la impedancia total tiene dos efectos fundamentales:
\begin{enumerate}
    \item \textbf{Reducción de la Corriente Total:} Por la ley de Ohm ($I_{total} = V_o / Z_{total\_media}$), una mayor impedancia implica que el motor en su conjunto extraerá menos corriente de la red eléctrica.
    \item \textbf{Caída de Tensión (Divisor de Tensión):} La bobina $L_3$ provoca una caída de tensión igual a $V_{L3} = I_{total} \times Z_{L3}$. En consecuencia, el voltaje útil que realmente llega al nodo $a$ (y que energiza los devanados del motor $L_1$ y $L_2$) será estrictamente menor que el voltaje original de la fuente $V_o$.
\end{enumerate}

\subsection*{Consecuencia Física: Reducción de la Velocidad}
Debido a que el voltaje efectivo en los devanados del motor es menor y fluye menos corriente a través de ellos, la intensidad del campo magnético rotatorio generado decae. Al tener un campo magnético más débil, el torque inducido en el rotor es menor, ocasionando que las aspas del ventilador giren a una frecuencia más baja. 

Esta configuración física representa la \textbf{velocidad media} del ventilador. Si conectáramos los nodos $d$ y $c$, sumaríamos en serie las impedancias de $L_3$ y $L_4$, lo cual aumentaría más la impedancia, reduciría más el voltaje en el nodo $a$, obteniendo finalmente la \textbf{velocidad baja}.

\subsection*{¿Por qué usar una bobina ($L_3$) y no una resistencia ($R_3$)?}
Es natural preguntarse si podríamos lograr el mismo efecto de reducir la velocidad colocando una resistencia $R_3$ en lugar de una bobina $L_3$. Matemáticamente, una resistencia también aumentaría la impedancia total ($Z_{total} = R_3 + Z_a$) y, por consiguiente, reduciría la corriente que llega al nodo $a$. 

Sin embargo, en la práctica \textbf{nunca se utilizan resistencias} para este propósito en motores de corriente alterna de este tamaño, debido a un factor físico fundamental: \textbf{la disipación de energía en forma de calor}.

\begin{itemize}
    \item \textbf{Comportamiento de una resistencia:} Las resistencias disipan \textit{potencia activa} ($P = I^2 R$). Para lograr la caída de tensión necesaria, la resistencia disiparía una cantidad enorme de energía térmica. Esto no solo representa un desperdicio altísimo de electricidad, sino que el calor acumulado derretiría los plásticos internos del ventilador, representando un grave riesgo de incendio.
    \item \textbf{Comportamiento de una bobina:} Los inductores son componentes \textit{reactivos}. Una bobina ideal no consume potencia activa ($P = 0$), sino que intercambia \textit{potencia reactiva} ($Q = I^2 X_L$) con la red eléctrica. Aunque las bobinas reales tienen una pequeña resistencia debido al alambre de cobre (lo que genera un ligero calentamiento), la pérdida térmica es minúscula comparada con la de una resistencia que genere la misma caída de tensión.
\end{itemize}

En resumen, la bobina se utiliza porque funciona como un \textbf{``limitador de corriente sin pérdidas de potencia activa''}, permitiendo variar la velocidad del ventilador de manera segura, fría y energéticamente eficiente.

\subsection*{Definiciones Adicionales: Potencia Activa y Reactiva}
Para comprender a cabalidad por qué las bobinas no generan calor excesivo a diferencia de las resistencias, es fundamental distinguir entre los tipos de potencia presentes en los circuitos de corriente alterna:

\begin{itemize}
    \item \textbf{Potencia Activa ($P$):} Es la potencia que realmente se consume y se transforma en trabajo útil o en calor disipado (como sucede en las resistencias). Se mide en \textbf{Vatios (W)} y representa la energía que la red eléctrica efectivamente entrega y no regresa al sistema. Matemáticamente se define como $P = V \cdot I \cdot \cos(\varphi)$, donde $\varphi$ es el ángulo de desfase entre la tensión y la corriente.
    
    \item \textbf{Potencia Reactiva ($Q$):} Es la potencia utilizada exclusivamente para crear y mantener los campos magnéticos en las bobinas o los campos eléctricos en los capacitores. Se mide en \textbf{Voltio-Amperios Reactivos (VAR)}. A diferencia de la potencia activa, esta energía no produce trabajo mecánico útil ni calor persistente; simplemente fluctúa u oscila, viajando de la fuente al componente reactivo y de regreso a la fuente en cada semiciclo de la onda alterna. Se define como $Q = V \cdot I \cdot \sin(\varphi)$.
\end{itemize}

\end{document}